\documentclass[10pt]{article}
\usepackage[margin=0.65in]{geometry}
\usepackage{graphicx}
\usepackage{booktabs}
\usepackage{array}
\usepackage{xcolor}
\usepackage{hyperref}
\usepackage{enumitem}
\hypersetup{colorlinks=true,linkcolor=blue!50!black,urlcolor=blue!50!black}
\graphicspath{{./}{reports/}{reports/figures/}{reports/figures/topology_zoo/}{reports/figures/toy_worlds/}{reports/figures/walk_geometry/}{docs/seminar/figures/}{docs/figures/2026-06-21-policy-storage/}}

\title{\textbf{HyMeKo Visual Atlas}\\[2pt]
\large Figures, plots, tables, and evidence map from the accessible workspace}
\author{Codex visual companion for Dr. Csaba Hajdu}
\date{2026-06-30}

\newcommand{\src}[1]{\par\footnotesize\textcolor{gray!70!black}{Source: \nolinkurl{#1}}\normalsize}
\newcommand{\note}[1]{\par\small\textit{#1}\normalsize}
\newcommand{\taggood}{\textcolor{green!45!black}{strong}}
\newcommand{\tagmixed}{\textcolor{orange!70!black}{mixed}}
\newcommand{\tagopen}{\textcolor{blue!60!black}{open}}
\newcommand{\tagneg}{\textcolor{red!70!black}{negative}}

\begin{document}
\maketitle

\begin{abstract}
\noindent This is the visual companion to
\nolinkurl{reports/2026-06-30-hymeko-exploration-summary.pdf}. It explains the
current HyMeKo state mostly through figures, plots, compact tables, and
evidence captions. It only uses artifacts accessible in this workspace: reports,
figures, benchmark CSVs, JSON summaries, and the current verification run.
\end{abstract}

\section{One-page evidence map}

This page is the reader's orientation key. It separates mature, publishable
evidence from open research bets, and it names the figure family that supports
each claim. The rest of the atlas expands each row into its figures, tables,
and missing discriminator.

\begin{table}[!htbp]\centering\small
\begin{tabular}{p{0.15\linewidth}p{0.28\linewidth}p{0.36\linewidth}p{0.11\linewidth}}
\toprule
Area & Main visual evidence & What it says & Readiness \\
\midrule
Cross-view IR & Cross-view grid and verification table & One HyMeKo source can
emit multiple concrete views whose recovered invariants agree. & \taggood \\
Performance substrate & Scaling and control-loop tables & Core emission,
expansion, and inference paths are fast enough to support tooling and demos. & \taggood \\
Topology matching & Topology map, zoo, arity plots & More edges do not
automatically help; controller topology/arity must match the plant. & \taggood \\
RL/control & Overnight RL, toy worlds, BC/off-policy plots & BC works; PPO or
off-policy refinement can collapse task success; fast toy lab is needed. & \tagopen \\
Vision/Soma & Readout ladder, spatial tree, RicciStim transfer & Position-aware
readout partially rehabilitates Soma vision; readout matters more than holonomy. & \tagmixed \\
FSR-LM & Architecture diagrams, bpb curve, speed curve & FSR pre-norm beats a
tiny transformer on bpb; top-k spike gate fixes dense long-context wall. & \tagopen \\
Holonomy & Probe table & Pure cycle holonomy is multiplicative transport, not
plain additive signed-walk aggregation. & \tagopen \\
\bottomrule
\end{tabular}
\caption{Compact map of the workspace's current evidence.}
\end{table}

\section{Cross-view consistency: the publishable kernel}

This is the strongest near-term contribution because it is a property of the
working system, not a training run. The verifier treats each emitted view as a
separate concrete syntax, extracts the same invariant, and checks agreement
against the IR-level query. That makes model drift a falsifiable engineering
failure rather than a vague documentation risk.

\begin{figure}[!htbp]\centering
\includegraphics[width=0.90\linewidth]{reports/figures/cross_view_consistency.png}
\caption{Machine-verified cross-view consistency. The key claim is a commuting
square over real emitted views: parse the concrete view back to an invariant,
and it agrees with the IR-level invariant.}
\src{reports/figures/cross_view_consistency.png; reports/2026-06-28-cross-view-consistency.md}
\end{figure}

\begin{table}[!htbp]\centering\small
\begin{tabular}{p{0.28\linewidth}p{0.24\linewidth}p{0.34\linewidth}}
\toprule
Domain & Views checked & Recovered invariant \\
\midrule
Robotics / kinematics & URDF, SDF, MJCF, DOT, Mermaid & Links, masses,
actuated joints, parent/child polarity, joint axes; exact for interchange views,
topological for diagram views. \\
Systems engineering & SysML requirements, DOT requirements & Requirements,
blocks, satisfy/allocate/derive relations, coverage, DAG/forest-style global
properties. \\
Drift demo & Single-source HyMeKo vs pairwise manual files & HyMeKo re-emission
keeps views synchronized; pairwise editing drifts by construction. \\
\bottomrule
\end{tabular}
\caption{The central formal-methods story is not a benchmark score; it is
machine-checkable consistency across concrete views.}
\end{table}

\section{Core scaling and control-loop substrate}

The infrastructure layer matters because the research claims sit on top of a
real compiler, transform registry, query engine, and tensor substrate. These
tables are not the paper's main novelty, but they show that the framework can
run the workflows that the papers and demos require.

\subsection{Evidence inventory}

\begin{table}[!htbp]\centering\small
\begin{tabular}{lrl}
\toprule
Artifact type under \nolinkurl{reports/} & Count & Meaning \\
\midrule
Markdown reports & 345 & Acceptance records and experiment narratives. \\
JSONL logs & 78 & Per-seed or per-run raw experiment traces. \\
PDF reports & 61 & Built technical reports, plans, and paper-style notes. \\
LaTeX sources & 56 & Rebuildable source for technical notes. \\
JSON summaries & 15 & Compact numeric outputs and benchmark summaries. \\
PNG figures at report root & 13 & Directly embedded summary plots. \\
\bottomrule
\end{tabular}
\caption{The repo is not just code; it is an evidence archive. This matters
because most claims can be traced to a report, figure, and raw-ish artifact.}
\src{PowerShell inventory over reports/}
\end{table}

\begin{table}[!htbp]\centering\small
\begin{tabular}{lrrrrrrrr}
\toprule
June report date & 15 & 16 & 17 & 18 & 19 & 20 & 21 & 29 \\
\midrule
Markdown reports & 17 & 7 & 10 & 12 & 14 & 14 & 13 & 12 \\
\bottomrule
\end{tabular}
\caption{June 2026 had several dense experiment days. The atlas focuses on the
highest-signal slices, not every report.}
\src{PowerShell grouping of reports/2026-06-*.md}
\end{table}

\begin{table}[!htbp]\centering\small
\begin{tabular}{lrrr}
\toprule
Corpus & Compile median ms & Star expansion median ms & MP forward median ms \\
\midrule
10 hyperedges & 0.514 & 0.000201 & 0.00684 \\
100 hyperedges & 3.074 & 0.00129 & 0.0580 \\
1000 hyperedges & 30.852 & 0.0128 & 0.554 \\
10000 hyperedges & 465.336 & 0.169 & 13.072 \\
\bottomrule
\end{tabular}
\caption{Selected scaling rows from the benchmark CSV. Expansion is tiny; message
passing stays practical into 10k-edge synthetic fixtures.}
\src{hymeko_bench/results/scaling.csv}
\end{table}

\begin{table}[!htbp]\centering\small
\begin{tabular}{llrr}
\toprule
Pipeline & Stage & Median us & P95 us \\
\midrule
Hypergraph & state assembly & 0.190 & 0.201 \\
Hypergraph & forward & 1.313 & 1.343 \\
Hypergraph & total & 1.593 & 1.624 \\
TD3 & total & 0.431 & 0.441 \\
PID+TD3 & total & 0.371 & 0.401 \\
\bottomrule
\end{tabular}
\caption{Older control-cycle benchmark: the hypergraph path has overhead versus
TD3/PID+TD3 but remains microsecond-scale in this harness.}
\src{hymeko_bench/results/control_cycle.csv}
\end{table}

\section{Hypergraph views and tensor intuition}

HyMeKo is most legible when the same object is viewed through several lenses:
star expansion, clique expansion, tensor/incidence views, and generated target
formats. These figures explain why the framework can talk simultaneously about
formal views, graph learning, and control. They also explain why arity and view
choice are not cosmetic: they change what information is preserved.

\begin{figure}[!htbp]\centering
\begin{minipage}{0.48\linewidth}\centering
\includegraphics[width=0.88\linewidth]{docs/seminar/figures/star_expansion.png}
\small Star expansion
\end{minipage}
\begin{minipage}{0.48\linewidth}\centering
\includegraphics[width=0.88\linewidth]{docs/seminar/figures/clique_expansion.png}
\small Clique expansion
\end{minipage}
\caption{Two canonical graph views of a hypergraph. Star expansion preserves
hyperedge identity via hub nodes; clique expansion collapses each hyperedge into
pairwise relations. This visual distinction underlies much of the arity and
view-consistency story.}
\src{docs/seminar/figures/star_expansion.png; docs/seminar/figures/clique_expansion.png}
\end{figure}

\begin{figure}[!htbp]\centering
\includegraphics[width=0.90\linewidth]{docs/seminar/figures/tensor_panels.png}
\caption{Tensor-view panels used in the seminar materials. These are the bridge
from declarative signed hypergraphs to tensorized learning/control substrates:
incidence, expansion, and message-passing views become concrete arrays.}
\src{docs/seminar/figures/tensor_panels.png}
\end{figure}

\begin{figure}[!htbp]\centering
\includegraphics[width=0.88\linewidth]{docs/seminar/figures/publications.png}
\caption{Publication/portfolio map from the seminar deck. It is useful as a
visual reminder that the project has multiple venue families, not one monolithic
paper.}
\src{docs/seminar/figures/publications.png}
\end{figure}

\section{Topology: the matching law}

The topology line asks a control-facing question: if the plant has a coupling
structure, should the controller mirror it? The supervised topology map says yes
at proxy scale: the matched controller wins every row, while the densest
complete controller does not act as a universal substitute. The important
claim is not ``this topology is best''; it is ``match beats capacity.''

\begin{figure}[!htbp]\centering
\includegraphics[width=0.93\linewidth]{reports/topology_map/topology_map.png}
\caption{Topology-to-performance matrix. The matched controller wins every
plant in this supervised structural proxy; the dense complete controller is not
a universal solution.}
\src{reports/topology_map/topology_map.png; reports/2026-06-27-topology-performance-map.md}
\end{figure}

\begin{figure}[!htbp]\centering
\includegraphics[width=0.96\linewidth]{reports/figures/topology_zoo/gallery.png}
\caption{Topology zoo: generic graphs, combinatorial designs, and extremal
high-symmetry constructions. These are candidate walk bases for structural
controllers.}
\src{reports/figures/topology_zoo/gallery.png; reports/2026-06-28-topology-zoo.tex}
\end{figure}

\begin{figure}[!htbp]\centering
\includegraphics[width=0.85\linewidth]{reports/figures/walk_geometry/coherence_bar.png}
\caption{Walk-basis frame coherence across topologies. Low coherence looks
mathematically attractive, but the companion control map shows that matching the
plant dominates any universal topology ranking.}
\src{reports/figures/walk_geometry/coherence_bar.png}
\end{figure}

\section{Arity: hypergraph lift is not free}

This section kills an attractive but wrong shortcut: lifting a graph plant into
a hypergraph does not automatically improve compute or accuracy. The lift helps
only when the plant itself is hypergraph-structured. That gives the project a
cleaner law: match arity to the plant, instead of treating hypergraphs as a free
upgrade over graphs.

\begin{figure}[!htbp]\centering
\begin{minipage}{0.48\linewidth}\centering
\includegraphics[width=\linewidth]{reports/figures/topology_zoo/compute_arity.png}
\small Compute cost
\end{minipage}
\begin{minipage}{0.48\linewidth}\centering
\includegraphics[width=\linewidth]{reports/figures/topology_zoo/accuracy_arity.png}
\small Accuracy on graph plants
\end{minipage}
\caption{The closed-neighbourhood hypergraph lift costs more and hurts graph
plants. The lift is a structural change, not a universal upgrade.}
\src{reports/figures/topology_zoo/compute_arity.png; reports/figures/topology_zoo/accuracy_arity.png}
\end{figure}

\begin{figure}[!htbp]\centering
\includegraphics[width=0.58\linewidth]{reports/figures/topology_zoo/arity_matching.png}
\caption{Arity matching. Graph controllers win on graph plants; hypergraph
controllers win on hypergraph plants. This is one of the cleanest control-side
principles in the repo.}
\src{reports/figures/topology_zoo/arity_matching.png; reports/2026-06-29-arity-study.tex}
\end{figure}

\section{RL/control: useful substrate, unresolved architecture win}

The RL part is an experimental engine rather than a finished victory claim. The
repo already has MuJoCo tasks, BC warm starts, PPO/off-policy refiners, GIF
rendering, checkpoints, and fast physics-free sanity worlds. The key lesson is
that task success and scalar return must be separated: a refinement run can make
the reward curve look better while the policy stops completing the task.

\begin{figure}[!htbp]\centering
\includegraphics[width=0.85\linewidth]{reports/figures/overnight_rl.png}
\caption{Overnight RL result. Behavior cloning produces useful pick-and-place
success; PPO refinement improves scalar return while actual lift/place success
collapses. This is reward misalignment, not yet an architecture win.}
\src{reports/figures/overnight_rl.png; reports/2026-06-29-rl-overnight.md}
\end{figure}

\begin{figure}[!htbp]\centering
\begin{minipage}{0.49\linewidth}\centering
\includegraphics[width=\linewidth]{reports/figures/toy_worlds/ablations.png}
\small Toy-world ablations
\end{minipage}
\begin{minipage}{0.49\linewidth}\centering
\includegraphics[width=\linewidth]{reports/figures/toy_worlds/nav_trajectories.png}
\small Navigation trajectories
\end{minipage}
\caption{Fast RL sanity worlds. These physics-free tests run real backbones in
seconds and are the right front-end before hour-scale MuJoCo runs.}
\src{reports/figures/toy_worlds/ablations.png; reports/figures/toy_worlds/nav_trajectories.png}
\end{figure}

\begin{figure}[!htbp]\centering
\begin{minipage}{0.32\linewidth}\centering
\includegraphics[width=\linewidth]{reports/figures/bc_breaks_wall.png}
\small BC break-through
\end{minipage}
\begin{minipage}{0.32\linewidth}\centering
\includegraphics[width=\linewidth]{reports/figures/offpolicy_collapse.png}
\small Off-policy collapse
\end{minipage}
\begin{minipage}{0.32\linewidth}\centering
\includegraphics[width=\linewidth]{reports/figures/fanuc_ppo_return.png}
\small Return curve
\end{minipage}
\caption{Manipulation learning diagnostics. The visual pattern is consistent:
imitation helps; naive refinement can move the scalar objective while damaging
task success.}
\src{reports/figures/bc_breaks_wall.png; reports/figures/offpolicy_collapse.png; reports/figures/fanuc_ppo_return.png}
\end{figure}

\begin{table}[!htbp]\centering\small
\begin{tabular}{lrrr}
\toprule
Fast RL arm & Reward / accuracy signal & Deploy ms & Params \\
\midrule
MLP & -0.157 & 0.12 & 3847 \\
SA-HSiKAN & -0.158 & 0.91 & 775 \\
HSiKAN-Chebyshev & -0.125 & 3.63 & 8071 \\
HSiKAN-CR & -0.133 & 4.67 & 8071 \\
Mixture & -0.146 & 5.17 & 11833 \\
\bottomrule
\end{tabular}
\caption{Fast bandit sanity results. SA-HSiKAN is far smaller/faster than vanilla
HSiKAN; HSiKAN-Chebyshev can improve reward but costs more rollout time.}
\src{reports/2026-06-29-fast-rl-sanity-suite.md}
\end{table}

\section{Vision/Soma: readout is the lever}

The vision line is best read as a diagnosis story. Early global-mean-pool
Soma/RicciStim looked falsified against simple controls. The later readout
sweeps show that the structure was often being compressed through the wrong
interface: discarding position erased the useful signal. Spatial-pyramid pooling
is the current reusable component because it keeps multi-scale position without
requiring a fixed flatten head.

\begin{figure}[!htbp]\centering
\includegraphics[width=0.82\linewidth]{reports/figures/soma_cheby_cell_readout_20260629.png}
\caption{Cell-by-readout sweep. The readout effect is much larger than the cell
effect: preserving position can flip Soma from apparent failure to useful signal.}
\src{reports/figures/soma_cheby_cell_readout_20260629.png; reports/2026-06-29-soma-cheby-cell-readout-sweep.md}
\end{figure}

\begin{figure}[!htbp]\centering
\begin{minipage}{0.49\linewidth}\centering
\includegraphics[width=\linewidth]{reports/figures/soma_posreadout_cluttered_20260629.png}
\small Single-vector readout plateau
\end{minipage}
\begin{minipage}{0.49\linewidth}\centering
\includegraphics[width=\linewidth]{reports/figures/soma_readout_ladder_20260629.png}
\small Readout ladder
\end{minipage}
\caption{Position-aware readout program. Mean pooling discards location;
attention helps; spatial pyramid/tree is the scalable middle; flatten is the
fixed-grid upper reference.}
\src{reports/figures/soma_posreadout_cluttered_20260629.png; reports/figures/soma_readout_ladder_20260629.png}
\end{figure}

\begin{figure}[!htbp]\centering
\begin{minipage}{0.49\linewidth}\centering
\includegraphics[width=\linewidth]{reports/figures/soma_ricci_tree_readout_20260630.png}
\small RicciStim transfer
\end{minipage}
\begin{minipage}{0.49\linewidth}\centering
\includegraphics[width=\linewidth]{reports/figures/soma_holonomy_group_ablation_20260629.png}
\small Holonomy group ablation
\end{minipage}
\caption{The visual verdict: position-aware spatial pooling helps RicciStim;
Z2/U(1) holonomy variants do not become the vision lever.}
\src{reports/figures/soma_ricci_tree_readout_20260630.png; reports/figures/soma_holonomy_group_ablation_20260629.png}
\end{figure}

\begin{table}[!htbp]\centering\small
\begin{tabular}{lrr}
\toprule
Readout & Parameters & Cluttered-MNIST accuracy \\
\midrule
Linear control & 23050 & 0.213 \\
Mean-pool & 2010 & 0.314 \\
Attention & 2027 & 0.466 \\
Spatial tree & 5227 & 0.537 \\
Flatten & 24890 & 0.617 \\
\bottomrule
\end{tabular}
\caption{Readout ladder as a table. The spatial-tree readout is the scalable
knee: about 87 percent of the attention-to-flatten gap at about one fifth of
flatten's parameters.}
\src{reports/2026-06-29-soma-spatial-tree-readout.md}
\end{table}

\section{FSR-LM: architecture, quality, and speed}

The FSR language-model line is the newest architecture bet. It is not a mature
scaling result; it is a tiny, controlled proof-of-concept that separates which
parts of the idea currently work. The pre-norm FSR mixer plus HSiKAN channel mix
beats a matched transformer on byte-level TinyShakespeare, while the spherical
Gomb residual loses. Sparse top-k fixes the long-context memory wall, but fused
attention still wins absolute throughput.

\begin{figure}[!htbp]\centering
\begin{minipage}{0.49\linewidth}\centering
\includegraphics[width=\linewidth]{reports/2026-06-29-fsr-lm-arch-model.png}
\small Model stack
\end{minipage}
\begin{minipage}{0.49\linewidth}\centering
\includegraphics[width=\linewidth]{reports/2026-06-29-fsr-lm-arch-mixer.png}
\small Fiber-Spike-Rotor mixer
\end{minipage}
\caption{FSR-LM architecture. The sequence mixer is signed, rotor-transported,
and spike-gated; the strongest reported variant uses standard pre-norm rather
than the Gomb spherical residual.}
\src{reports/2026-06-29-fsr-lm-arch-model.png; reports/2026-06-29-fsr-lm-arch-mixer.png}
\end{figure}

\begin{figure}[!htbp]\centering
\begin{minipage}{0.49\linewidth}\centering
\includegraphics[width=\linewidth]{reports/2026-06-29-fsr-lm-phase1.png}
\small TinyShakespeare A/B
\end{minipage}
\begin{minipage}{0.49\linewidth}\centering
\includegraphics[width=\linewidth]{reports/2026-06-29-fsr-lm-speed.png}
\small Throughput vs sequence length
\end{minipage}
\caption{FSR quality and speed. FSR pre-norm beats the matched transformer on
tiny byte-LM bpb; hard top-k sparsity prevents the dense long-context wall, but
absolute speed still trails fused transformer kernels.}
\src{reports/2026-06-29-fsr-lm-phase1.png; reports/2026-06-29-fsr-lm-speed.png}
\end{figure}

\begin{table}[!htbp]\centering\small
\begin{tabular}{lrrr}
\toprule
Model & Val bpb median & Tokens/s median & Params \\
\midrule
FSR sphere & 2.7702 & 11780 & 367328 \\
FSR pre-norm & 2.4791 & 11161 & 368960 \\
Transformer & 2.5698 & 35080 & 369324 \\
\bottomrule
\end{tabular}
\caption{FSR Phase 1 result table. The win belongs to FSR pre-norm; the sphere
residual loses at full budget.}
\src{reports/2026-06-29-fsr-lm-phase1-ab.json}
\end{table}

\begin{table}[!htbp]\centering\small
\begin{tabular}{rrrr}
\toprule
Sequence length & Dense FSR tok/s & FSR k=16 tok/s & Transformer tok/s \\
\midrule
128 & 5510 & 5795 & 20789 \\
256 & 4527 & 11365 & 35275 \\
512 & 583 & 14765 & 46808 \\
1024 & OOM & 11138 & 106329 \\
\bottomrule
\end{tabular}
\caption{Top-k sparse FSR changes the scaling regime: dense collapses/OOMs while
k=16 remains feasible.}
\src{reports/2026-06-29-fsr-lm-throughput-vs-T.json}
\end{table}

\section{Holonomy: multiplicative transport is the signal}

This section is deliberately small but important. It decides what kind of
``structure'' the gauge line needs. A cycle-parity label is not an additive
aggregate; it is a product accumulated around a loop. The current probe shows
that a multiplicative transport reader solves the task exactly, while an
additive signed-walk reader, a flat MLP, and a linear guard stay near chance.
That supports rotor/transport-style mechanisms for holonomy-heavy tasks and
warns against overclaiming plain additive HSiKAN on pure parity.

\begin{table}[!htbp]\centering\small
\begin{tabular}{lrrp{0.34\linewidth}}
\toprule
Arm & Accuracy & Params & Meaning \\
\midrule
Transport / rotor-like reader & 1.0000 & 49 & Reads the product around the
cycle exactly. \\
Additive signed-walk reader & 0.4805 & 225 & The additive HSiKAN-like
mechanism does not expose pure parity holonomy here. \\
MLP & 0.5913 & 1505 & Near chance; parity is hard from raw local signs. \\
Linear guard & 0.5054 & 13 & Confirms no local linear leakage. \\
\bottomrule
\end{tabular}
\caption{Current holonomy probe run performed during the exploration. This is a
decisive visual table: pure cycle holonomy is multiplicative transport, not a
plain additive signed-walk sum.}
\src{Command: .\\.venv\\Scripts\\python.exe -m hymeko_rl.holonomy_probe --ring-size 12 --seeds 2}
\end{table}

\begin{table}[!htbp]\centering\small
\begin{tabular}{p{0.22\linewidth}p{0.31\linewidth}p{0.31\linewidth}}
\toprule
Question & What the probe answers & What it does not yet answer \\
\midrule
Is parity locally visible? & No. Linear guard stays at chance, so the label is
not leaking through a local linear cue. & It does not test every possible
feature-engineered baseline. \\
Does additive message passing read pure holonomy? & No. The additive reader is
near chance on this task. & It does not falsify HSiKAN on non-parity tasks. \\
Does transport read holonomy? & Yes. Multiplicative transport reaches 1.0
accuracy in the current run. & Continuous Spin holonomy and aromaticity variants
still need full execution. \\
\bottomrule
\end{tabular}
\caption{Interpretation guard for the holonomy probe. The result is sharp, but
its scope is parity-like cycle holonomy, not all structural learning.}
\end{table}

\section{Signed-link and rotor line}

The signed-link line contains some of the strongest-looking numbers in the repo,
but it also contains the strongest warnings about leakage and attribution. The
safe framing is: strict protocol design and structural features are valuable;
rotor embeddings are parameter-efficient and inductive; the rotor geometry
itself is not yet the measured source of AUROC gains. This section therefore
prioritizes audit discipline over headline SOTA language.

\begin{figure}[!htbp]\centering
\begin{minipage}{0.49\linewidth}\centering
\includegraphics[width=\linewidth]{docs/seminar/figures/rotor_param_efficiency.png}
\small Parameter efficiency
\end{minipage}
\begin{minipage}{0.49\linewidth}\centering
\includegraphics[width=\linewidth]{docs/seminar/figures/rotor_pareto.png}
\small Pareto view
\end{minipage}
\caption{Cayley-rotor signed-link evidence. The honest reading from the reports:
rotor embeddings are parameter-efficient and inductive, but later ablations say
the rotor geometry itself is not the sole load-bearing ingredient for AUROC. The
value is structural features plus strict leakage discipline.}
\src{docs/seminar/figures/rotor_param_efficiency.png; docs/seminar/figures/rotor_pareto.png; ROADMAP.md; DONE.md}
\end{figure}

\begin{table}[!htbp]\centering\small
\begin{tabular}{lrrrp{0.34\linewidth}}
\toprule
Model/result family & Alpha & OTC & Slashdot/Epinions & Interpretation \\
\midrule
Gomb strict tuned & 0.897 & 0.915 & Slashdot 0.902 / Epinions 0.943 &
Honest strict baseline; avoids the sigma-leakage that inflated earlier cycle
numbers. \\
Optuna-best HSiKAN & 0.996 & 0.993 & -- & Strong iso-protocol numbers, but
reported as transductive/leakage-sensitive and must be framed carefully. \\
Rotor-vs-MLP embed & tied & tied & tied & Foundational ablation: rotor
geometry is not load-bearing on AUROC; structural features are the safer claim. \\
\bottomrule
\end{tabular}
\caption{Signed-link takeaway table. The graph-learning line is valuable, but
the publishable story is strict protocols and structural priors, not unqualified
SOTA claims.}
\src{docs/SOTA_RESULTS.md; reports/2026-06-18-rotor-vs-mlp-embed-ablation.md}
\end{table}

\begin{table}[!htbp]\centering\small
\begin{tabular}{p{0.24\linewidth}p{0.30\linewidth}p{0.30\linewidth}}
\toprule
Risk & Why it matters & Current mitigation \\
\midrule
Sigma-product leakage & Cycle products can silently include test-edge sign
information under transductive conventions. & Shuffle gates and strict
train-edge enumeration identify inflated numbers. \\
Geometry overclaim & Rotor representations can look novel even when an MLP
embedding ties their AUROC. & Rotor-vs-MLP ablation reframes the win as
structural features / inductive compactness. \\
Protocol mismatch & Published baselines may use different dedup or reachability
rules. & Protocol-matched SiGAT comparison and planned tier-2/3 audits. \\
\bottomrule
\end{tabular}
\caption{Why the signed-link section should be read cautiously. The contribution
is stronger when it leads with protocol honesty and parameter efficiency.}
\end{table}

\section{Policy-as-hypergraph storage}

A useful bridge between the framework and RL is that trained policies can become
HyMeKo artifacts themselves. The point is not only storing weights compactly; it
is making policy provenance, reconstruction, and rendering part of the same
hypergraph workflow. This is currently proven on cart-pole-scale policy storage
and should be extended to manipulation policies once BC policies are selected as
the reliable demo baseline.

\begin{figure}[!htbp]\centering
\begin{minipage}{0.32\linewidth}\centering
\includegraphics[width=\linewidth]{docs/figures/2026-06-21-policy-storage/01-weight-matrix-star-expansion.png}
\small Weight matrix as star expansion
\end{minipage}
\begin{minipage}{0.32\linewidth}\centering
\includegraphics[width=\linewidth]{docs/figures/2026-06-21-policy-storage/03-actor-critic-dataflow.png}
\small Actor-critic dataflow
\end{minipage}
\begin{minipage}{0.32\linewidth}\centering
\includegraphics[width=\linewidth]{docs/figures/2026-06-21-policy-storage/06-roundtrip-verification.png}
\small Round-trip verification
\end{minipage}
\caption{Policy storage as HyMeKo hypergraph. A trained policy can be serialized
as a valid HyMeKo graph, reconstructed, and checked bit-exact in the reported
cart-pole storage work. This is a bridge between framework and RL: models become
first-class graph artifacts.}
\src{docs/figures/2026-06-21-policy-storage/; reports/2026-06-21-policy-weight-storage.md}
\end{figure}

\begin{table}[!htbp]\centering\small
\begin{tabular}{lrrp{0.38\linewidth}}
\toprule
Storage tier & HyMeKo size & External blob & Meaning \\
\midrule
T0 decimal & 332 KB & none & Human-readable but large. \\
T1 base64/auto & 142--148 KB & none & More compact textual embedding. \\
T2 content-addressed & 9 KB & NPZ blob & Small HyMeKo manifest plus verified
external tensor payload. \\
\bottomrule
\end{tabular}
\caption{Policy-storage tiers reported on 2026-06-21. The important idea is not
compression alone; it is provenance and structural round-trip as part of the
model artifact.}
\src{DONE.md, 2026-06-21 policy-storage entries}
\end{table}

\section{Thirteen pursuit lanes: what is still missing}

The atlas is evidence-heavy, but evidence is only useful if it points to the
next decisive measurement. This board lists the thirteen lanes that would
convert the repository's current state into a set of finished claims. Each lane
has an existing artifact, a missing discriminator, and a smallest credible next
action.

\begin{table}[!htbp]\centering\scriptsize
\begin{tabular}{rp{0.19\linewidth}p{0.27\linewidth}p{0.27\linewidth}p{0.14\linewidth}}
\toprule
\# & Lane & What exists now & Missing discriminator & Near action \\
\midrule
1 & Cross-view paper & Real CLI emitters, robotics+SysML witnesses, drift demo,
proof bundle & Reviewer-proof positioning and final article polish & Submit
T-SMC/SMC-style package. \\
2 & CI consistency gate & 22-test cross-view suite, extractor strategies & CI
integration and optional Gazebo/torch-dataflow extensions & Add verifier as
regression gate. \\
3 & Topology matching & 8/8 supervised topology map, exact equivariance guard &
Closed-loop correlation on real control & Run matched-vs-mismatched controller
on cart-pole/arm. \\
4 & Arity matching & Graph-vs-hypergraph plant reversal proves matching law &
More plant families and task types & Repeat on contact/gait hypergraph plants. \\
5 & RL refinement & BC works; PPO/off-policy can collapse success & BC-anchor
and success-aligned reward proving success does not collapse & Make BC anchor
default for refine sweeps. \\
6 & Fast RL lab & Bandit/grid/hex/collab toys; prescreen module & Correlation
between toy ranking and MuJoCo ranking & One real-scale confirm per toy-chosen
winner. \\
7 & Holonomy & Transport solves parity; additive/MLP near chance & Continuous
Spin holonomy and aromaticity versions through production arms & Extend probe
beyond T1 parity. \\
8 & StructuralActor & B-power collapse, speed/parameter win, Steiner basis &
Closed-loop actuator use, not just supervised target fit & Wire to a control
task. \\
9 & FSR-LM & TinyShakespeare bpb win; sparse top-k long-context feasibility &
Scale beyond 1 MB corpus, ablations of sign/rotor/spike/cell & Enwik8/text8
subset with spike\_k=16. \\
10 & Soma readout & Spatial tree/pyramid improves fixed and variable anchors &
Natural-image or cortical-predictivity validation & Export
SpatialPyramidPool and test in HyMeYOLO/RL vision. \\
11 & Signed-link & Strict Gomb baseline, leakage audit, rotor Pareto evidence &
Final strict baseline table and no-overclaim narrative & Finish tier-2/3 and
shuffle audits. \\
12 & Policy storage & Bit-exact policy round-trip and compact T2 storage &
Generalize from cart-pole to manipulation policies & Store/render pick-place
BC policies as HyMeKo artifacts. \\
13 & Demo narrative & Many figures, GIFs, web views, editor surfaces & One
curated path that a collaborator can run end-to-end & Script ``one source,
many views, one policy'' demo. \\
\bottomrule
\end{tabular}
\caption{The missing-detail board. These are the thirteen lanes that convert the
atlas from evidence inventory into an execution plan.}
\end{table}

\begin{table}[!htbp]\centering\small
\begin{tabular}{p{0.22\linewidth}p{0.26\linewidth}p{0.25\linewidth}p{0.17\linewidth}}
\toprule
Scenario / artifact & Current accessible outputs & Main risk visible in atlas & Best next measurement \\
\midrule
Pick-and-place & BC/PPO checkpoints, overnight plot, GIF folders & PPO improves
return but collapses success & BC-anchor A/B on lift/place success. \\
Coin-toss / Galambos & Collaborative GIFs, checkpoints, delivery table & CTDE
undertrained; single agent currently wins & Longer CTDE refine with anchor. \\
Planar walker & Walker checkpoint and GIF & Partial forward movement, not full
goal reaching & 300-iter or shorter-goal confirm. \\
Cart-pole storage & Policy graph, storage figures, GIFs & Easy task; structure
not load-bearing & Use as storage/provenance demo, not architecture proof. \\
FSR-LM & Phase0/1 plots, architecture diagrams, speed plots & Tiny scale and
absolute speed gap & Larger corpus plus fused-kernel decision. \\
Soma/RicciStim & Readout ladders and Ricci tree figures & Still below flatten
and conventional baselines & Spatial pyramid in detection/RL vision. \\
Cross-view IR & Proof PDF, cross-view figure, JSON verification & Claim must stay
about consistency, not raw generator speed & Add CI gate and final paper polish. \\
\bottomrule
\end{tabular}
\caption{Scenario matrix: what is runnable or visible now, why it is not yet a
finished claim, and the smallest next measurement that would change the state.}
\end{table}

\begin{table}[!htbp]\centering\small
\begin{tabular}{lrrp{0.42\linewidth}}
\toprule
Artifact family & Count & Where & Why it matters \\
\midrule
Checkpoints & 95 files & \nolinkurl{checkpoints/} & The learning lines are not
only reports; there are stored policies/models for replay, rendering, and
ablation. \\
Rendered GIFs & 91 files & \nolinkurl{reports/gifs/} & The RL/control claims can
be inspected behaviorally, not only through scalar metrics. \\
Plans & 100+ dirs & \nolinkurl{docs/plans/} & The project has an explicit
plan-first discipline; unfinished work should be judged against these gates. \\
Reports & 345 md & \nolinkurl{reports/} & Most completed tasks have an
acceptance note; the atlas samples the strongest, not all. \\
\bottomrule
\end{tabular}
\caption{Accessible artifact inventory beyond the embedded figures. The missing
work is mostly not lack of code; it is choosing the decisive measurement.}
\end{table}

\section{Publication and execution matrix}

This matrix turns the figures into venue strategy. Different communities should
see different leads: formal methods should see cross-view consistency; control
should see topology and arity matching; ML architecture should see FSR plus its
ablations; robotics collaborators should see BC, scenarios, and demo readiness.
Keeping those leads separate prevents one weak line from diluting a stronger
one.

\begin{table}[!htbp]\centering\scriptsize
\begin{tabular}{p{0.18\linewidth}p{0.30\linewidth}p{0.28\linewidth}p{0.16\linewidth}}
\toprule
Venue family & Lead claim & Evidence to show visually & Blocking next step \\
\midrule
IEEE T-SMC / SMC & Canonical hypergraph IR with machine-verified cross-view
consistency & Cross-view grid, systems-engineering table, drift demo, proof
bundle & Polish article and supplementary; keep claims centered on consistency. \\
Formal methods / MDE & Model consistency becomes a machine-checked property,
not a convention & Commuting-square diagram, extractor table, Z3 sat/unsat
summary & Position against EMF/MLIR/global model management. \\
Control / CPS & Matched topology and matched arity govern controller
representation & Topology heatmap, arity matching plot, StructuralActor speed
table & Closed-loop correlation on real control tasks. \\
RL robotics collaboration & HyMeKo supplies declarative structural RL scenarios
and fast pre-screening & Overnight RL figure, toy-world ablations, BC collapse
plots & Repair PPO refinement with BC anchor and success-aligned reward. \\
Graph ML / signed networks & Strict leakage-audited structural features are a
methodological contribution & SOTA tables, shuffle gates, transfer grid,
rotor-vs-MLP ablation & Finish strict baseline audit and avoid rotor-overclaim. \\
ML architecture workshop & FSR mixer is a plausible attention alternative at
tiny scale & FSR architecture, bpb plot, sparse speed plot & Scale beyond
TinyShakespeare; ablate sign/rotor/spike/channel mix. \\
Vision / neuro & Position-aware readout revives part of Soma/RicciStim, but
holonomy is not the vision lever & Readout ladder, Ricci tree transfer, holonomy
group ablation & Try stronger scalable readouts and brain-predictivity metrics. \\
\bottomrule
\end{tabular}
\caption{Where each evidence block can go, and what is still missing.}
\end{table}

\section{Visual conclusion}

The shortest path through the atlas is not chronological. It is a set of
decision figures: one for the publishable framework kernel, one for the
structure-matching law, one for the RL warning, one for readout, one for the LM
bet, and one for holonomy. These are the figures to use when a collaborator has
only a few minutes.

\begin{table}[!htbp]\centering\small
\begin{tabular}{p{0.30\linewidth}p{0.55\linewidth}}
\toprule
If you show one figure & Use the cross-view consistency figure: it is the
publishable kernel. \\
If you show one control result & Use topology map plus arity matching: structure
helps when it matches the plant. \\
If you show one RL warning & Use overnight RL: scalar return improvement can
hide success collapse. \\
If you show one vision result & Use the readout ladder: position-aware pooling
is the lever. \\
If you show one new architecture result & Use FSR bpb plus speed: quality is
promising, speed needs fused kernels. \\
If you show one holonomy result & Use the transport/additive/MLP table:
holonomy is multiplicative. \\
\bottomrule
\end{tabular}
\caption{The shortest visual route through the entire workspace.}
\end{table}

\end{document}
